\documentclass[12pt,a4paper]{report}
\usepackage[dvips]{graphicx}           
\usepackage[portuges]{babel}           
\usepackage{color}            
\usepackage[latin1]{inputenc}           
\usepackage[normalem]{ulem}
\setlength{\textwidth}{160mm}
\setlength{\textheight}{220mm}
\setlength{\topmargin}{-8pt}
\setlength{\oddsidemargin}{0mm}
\usepackage{amsmath,amsfonts, amssymb}
\usepackage{enumitem}% http://ctan.org/pkg/enumitem

\begin{document}

  \begin{minipage}[b]{0.1 \linewidth}
     \includegraphics[width=.9cm]{ufs_vertical_positiva.png}
  \end{minipage}\hfill	
  \begin{minipage}[b]{0.85 \linewidth}
	       \textsc{\large Universidade Federal de Sergipe}\\
	       \textsc{\large Departamento de Estat�stica e Ci�ncias Atuariais}\\
	       Disciplina: ESTAT0078 -- Infer�ncia I\\
	       Professor: Sadraque E. F. Lucena
\end{minipage}

\vspace{0.8cm}
\begin{center} {\bf Lista de Exerc�cios 10} \end{center}
\vspace{0.8cm}

	\begin{enumerate}[label=10.\arabic*)]
		
		% Quest�o 1
		\item
		Seja $X_1,\ldots,X_n$ uma amostra aleat�ria de uma vari�vel aleat�ria $X$. Em cada item, determine uma estat�stica suficiente para o par�metro indicado, utilizando o Crit�rio da Fatora��o de Neyman. Apresente a fun��o de verossimilhan�a e identifique explicitamente os fatores $h(\underset{\widetilde{\phantom{a}}}{x})$ e $g_\theta(T(\underset{\widetilde{\phantom{a}}}{x}))$.
		\begin{enumerate}[label=\alph*)]
			\item $X\sim\operatorname{Binomial}(2,\theta)$, com $0<\theta<1$.
			\item $X$ tem densidade $f(x;\theta)=\theta x^{\theta-1},\qquad 0<x<1,\quad \theta>0$.
			\item $X\sim\operatorname{Exponencial}(\theta)$, com densidade $f(x;\theta)=\frac{1}{\theta}e^{-x/\theta},\qquad x>0,\quad \theta>0$.
			\item $X\sim N(0,\sigma^2)$, em que $\sigma^2>0$ � desconhecido.
			\item $P(X=x;\theta)=\theta(1-\theta)^{x-1},\qquad
			x=1,2,\ldots,\quad 0<\theta<1$.
		\end{enumerate}
		
	\end{enumerate}
	
	\newpage
	
	\noindent Respostas:
	
	\begin{enumerate}[label=10.\arabic*)]
		
		% Quest�o 1
		\item 
		\begin{enumerate}
			\item $T(\underset{\widetilde{\phantom{a}}}{x}) = \sum\limits_{i=1}^n X_i$\\
			$h(\underset{\widetilde{\phantom{a}}}{x}) = \prod\limits_{i=1}^n \binom{2}{x_i}$\\
			$g_\theta(T(\underset{\widetilde{\phantom{a}}}{x})) = \theta^{\sum\limits_{i=1}^n X_i} (1-\theta)^{2n - \sum\limits_{i=1}^n X_i}$
			
			\item $T(\underset{\widetilde{\phantom{a}}}{x}) = \prod\limits_{i=1}^n X_i$\\
			$h(\underset{\widetilde{\phantom{a}}}{x}) = 1$\\
			$g_\theta(T(\underset{\widetilde{\phantom{a}}}{x})) = \theta^n \left( \prod\limits_{i=1}^n X_i \right)^{\theta-1}$
			
			\item $T(\underset{\widetilde{\phantom{a}}}{x}) = \sum\limits_{i=1}^n X_i$\\
			$h(\underset{\widetilde{\phantom{a}}}{x}) = 1$\\
			$g_\theta(T(\underset{\widetilde{\phantom{a}}}{x})) = \frac{1}{\theta^n} \exp\left( -\frac{\sum\limits_{i=1}^n X_i}{\theta} \right)$
			
			\item $T(\underset{\widetilde{\phantom{a}}}{x}) = \sum\limits_{i=1}^n X_i^2$\\
			$h(\underset{\widetilde{\phantom{a}}}{x}) = (2\pi)^{-n/2}$\\
			$g_{\sigma^2}(T(\underset{\widetilde{\phantom{a}}}{x})) = (\sigma^2)^{-n/2} \exp\left( -\frac{\sum\limits_{i=1}^n X_i^2}{2\sigma^2} \right)$
			
			\item $T(\underset{\widetilde{\phantom{a}}}{x}) = \sum\limits_{i=1}^n x_i$\\
			$h(\underset{\widetilde{\phantom{a}}}{x}) = 1$\\
			$g_\theta(T(\underset{\widetilde{\phantom{a}}}{x})) = \theta^n (1-\theta)^{\sum\limits_{i=1}^n x_i - n}$\
		\end{enumerate}
		
	\end{enumerate}
	
\end{document}